\section{Reproducibility}
\label{sec:repro}

This appendix summarizes the main steps used to reproduce the simulations, validation tables, and figures reported in the manuscript.

\subsection{Li et al.\ (2023) Table~2 cases}
\begin{verbatim}
# Validate all 11 published cases (reads summary.json)
python3 scripts/validate_li2023_all.py --cases all \
  --tag v9_table2_extended

# Run all 11 FEM cases (slow; ~30+ min each)
python3 scripts/run_li2023_cases.py --cases all \
  --tag v9_table2_extended

# Run calibration subset only (cases 1, 6, 87, 90)
python3 scripts/run_li2023_cases.py --cases calibration \
  --tag v8_full_four_cases_lowE_single
\end{verbatim}

Data definitions: \repoPath{scripts/li2023\_data.py}, \repoPath{config/data/li2023\_tables.json}.

\begin{verbatim}
# Regenerate long-format validation CSV and publication plots
.venv_lineheating/bin/python \
  scripts/make_publication_validation_figures.py

# Regenerate journal-style FEM surface panels
.venv_lineheating/bin/python \
  scripts/make_publication_surface_panels.py

# Execute calibration, thermal-unit audit, and verification scaffold
.venv_lineheating/bin/python \
  scripts/execute_planned_work_items.py
\end{verbatim}

These commands write long-format validation CSV files, metrics JSON files, validation plots, journal-style FEM surface panels, velocity-dependent calibration outputs, keel thermal-unit audit tables, and verification/energy scaffold files under \repoPath{results/} and the publication figure folders.

\subsection{Induction-style line-source process comparison}
\begin{verbatim}
# Collapsed full-line induction source with quench (all 11 cases)
python3 scripts/run_li2023_induction_cases.py --cases all \
  --tag v10_induction_simul_quench \
  --heat-mode simultaneous --skip-existing

# Sequential line-source induction with quench (multi-pass subset)
python3 scripts/run_li2023_induction_cases.py --cases 6,7,87,88,89 \
  --tag v11_induction_seq_quench \
  --heat-mode sequential --skip-existing

# Volumetric induction skin-depth source
python3 scripts/run_li2023_induction_cases.py --cases 6,7,87,88,89 \
  --tag v12_induction_skin_seq_quench \
  --heat-mode sequential --heat-source-mode induction_skin \
  --induction-frequency 10000 --induction-relative-permeability 100 \
  --induction-electrical-resistivity 2.5e-7 --induction-efficiency 0.8 \
  --skip-existing

# Validate available v11 summaries
python3 scripts/validate_li2023_all.py --cases 6,7,87,88,89 \
  --tag v11_induction_seq_quench
\end{verbatim}

The v10 configuration uses a single full-line heating event for multi-pass rows by reducing velocity to $v/N$. The v11 configuration keeps the same line heat source and quench model but repeats the pass sequentially. The v12 command switches from prescribed surface band heating to the \code{induction\_skin} volumetric source.

\subsection{Keel-plate cases}
\begin{verbatim}
python scripts/cases/run_keel_fast_parallel.py
python scripts/cases/run_keel_calibrated.py
python scripts/cases/run_keel_enhanced.py
python scripts/cases/run_keel_li2023_scaled.py --run
\end{verbatim}

\subsection{Build manuscript}
\begin{verbatim}
cd publication2
latexmk -pdf main.tex
\end{verbatim}
Figures are read from the self-contained \repoPath{publication2/figures/} folder.
